\scanpage{63}
$\zeta z$ est une fonction impaire, \deleted{\unsure{non périodique}. On a} \inserted{admettant les \unsure{nbs} $\omega$ pour pole simple (avec 1 comme résidu)}
\[
\zeta(z+2\omega_1)-\zeta z=c^{\mathrm{te}}=2\zeta\omega_1=2\eta_1\qquad
\text{on définit de }\unsure{\overline{\mathrm m}|=}\ \eta_2\text{ et }\eta_3
\]

On établit par une intégration de $\int z\,\unsure{p|\wp}z\,dz$ sur un contour convenable
\[
\omega_2\eta_1-\omega_1\eta_2=\frac{\pi i}{2}\qquad\text{etc. ; d'ailleurs }\quad\eta_1+\eta_2+\eta_3=0
\]

La fonction $\zeta z$ permet de calculer l'intégrale $\displaystyle\int\frac{Z\,dZ}{\sqrt{4Z^3-g_2Z-g_3}}$ on pose $Z=pz$ d'où $\displaystyle\int=\int pz\,dz=-\zeta z+c^{\mathrm{te}}$

\deleted{D'une manière \unsure{précise|générale},}

\uline{Formule de décomposition de Hermite}. Soit $f(z)$ une fonction elliptique de périodes $2\omega_1$ et $2\omega_2$ ; $a,b,\mathord{--}$ ses poles, distincts à une période près deux à deux,
\[
\frac{A_\alpha}{(z-a)^\alpha}+\mathord{--}+\frac{A_1}{z-a}\ ;\qquad
\frac{B_\beta}{(z-b)^\beta}\mathord{--}+\frac{\unsure{b|B}_{\unsure{1|2}}}{z-b}\ ;\ \mathord{---}
\]
les parties principales de ses poles ; On a alors
\[
\begin{aligned}
f(z)=c^{\mathrm{te}}+{}&\left\{A_1\zeta(z-a)+A_2p(z-a)-\frac{A_3}{2!}p'(z-a)+\mathord{--}+(-1)^\alpha\frac{A_\alpha}{(\alpha-1)}p^{(\alpha-2)}(z-a)\right\}\\
+{}&\{B_1\zeta(z-b)+\mathord{---}\qquad\qquad\}\\
&\mathord{-}\quad\mathord{-}\quad\mathord{-}\quad\mathord{-}\quad\mathord{-}\quad\mathord{-}\quad\mathord{-}
\end{aligned}
\]

\fbox{\uline{fonction $\sigma z$}} Nous savons intégrer toutes les fonctions elliptiques au moyen de fonctions connues, si nous savons intégrer $\zeta z$. Or, son intégrale est complexe, \deleted{et} admet \unsure{tous} les points $\omega$ pour points logarithmiques. Cela suggère de chercher $\sigma z$ tel que $\int\zeta z\,dz=\log\sigma z$\qquad $\sigma z$ s'annulant aux points $\omega$\unsure{--}

% archive sheet number: 62
